Property-based testing (PBT) is a widely used automated testing technique that shifts the burden of writing individual test cases from the user to the automated tool. Unlike traditional unit testing, which relies on manually crafted examples, PBT uses test input generators to automatically produce a family of test cases, focusing on potential bugs that violate specified properties. PBT originated in the functional programming community, where programs are often pure and test cases can be generated simply as input arguments. For example, to test the merge function in a merge sort algorithm, a generator might produce only sorted lists as inputs. 
Over time, PBT has been extended to more complex and stateful systems, where test cases are sequences of API calls to libraries, user actions in GUIs~\cite{Quickstrom}, or request messages in distributed systems.
Despite these differences in test input forms, the core principle of PBT remains the same: it reduces untargeted and random exploration by using properties, written in expressive specification languages, to guide and specialize test input generation before actual testing occurs.

However, directly applying the core PBT principle becomes challenging in systems where certain APIs provided by system under test (SUT) do not immediately change the system state. For example, consider a storage system that handles read and write request messages from users. When a request is sent, the stored value may not be updated immediately due to network latency. In such cases, test inputs—such as read and write requests—cannot fully guide the exploration during testing, because the order in which messages are delivered and processed is an effect outside the control of the test input generator.
We call these \emph{latent effects}: hidden behaviors that are not directly observable or controllable through the API provided by the SUT, but can change the system state later. Existing PBT frameworks cannot see or control latent effects, even though they are often the source of critical bugs—such as those caused by message reordering in distributed systems. As a result, PBT is less effective for systems with latent effects, since test input generation alone cannot fully guide or constrain the exploration of possible executions.

In order to design a PBT framework that is aware of latent effects, there are two immediate challenges that need to be overcome. First, we need to provide specifications expressive enough to capture the properties of latent effects that cause critical bugs within the SUT. Consider the storage system example: handling a request from the client side may trigger a lot of message passing between different components of the distributed system that are not exposed by the storage system's APIs (i.e., the read and write requests). The user should be able to specify these latent messages within the specification language.  
Second, we need some way to leverage these specifications to control latent effects and appropriately bias our search procedure towards executions that are likely to reveal violations of desired behaviors. Unlike traditional PBT, where generators only need to provide test inputs satisfying the precondition of the property to be tested, the desired generator should guide the exploration over latent messages in specific orders. Requiring the user to come up with a global order of latent messages that may trigger bugs is not practical; thus, smarter automated techniques are needed to facilitate test input generation.

In this paper, we present a unified solution to both of these challenges. Building on the trace-based approach common in PBT frameworks for stateful systems, we express properties as global predicates over traces of API invocations. Crucially, we make latent actions explicit by modeling them as latent events within these traces. Latent effects are abstracted as capabilities—representing the potential for certain events to occur in the future.
To capture these behaviors, we extend recent trace-based refinement type systems with automata that describe programs with opaque internal state~\cite{ZYDJ24}, introducing a new ``prophecy automata'' component that specifies latent effects that may arise in the future. Our test generator is thus empowered not only to generate API invocations for the SUT, but also to schedule both observable and latent events, controlling their order to systematically explore possible executions.
Rather than requiring users to specify a global trace of (latent) events, our approach allows users to provide only local relationships between events—for example, specifying that a write request should eventually trigger a write response. More complex behaviors, such as the possibility that ``a read immediately after a write may trigger an inconsistency bug,'' are automatically discovered by our test generator synthesizer, which composes local specifications and global properties to guide the search for violations.